\begin{figure*}[t]
\centering
\begin{scriptsize}
\begin{subfigure}{0.47\linewidth}
\centering
\begin{tikzpicture}
\tkzKiviatDiagram[scale=0.35,
        label distance=.5cm,
        radial  = 10,
        gap     = 1,
        lattice = 5,
        step=.05]
        {Trash\\pile,Fire,Car,Trash\\pile,Person,Building,Campsite,Forest\\fire,Road\\construction,Crowd}
\tkzKiviatLine[very thick,color=color0](19.570,41.300,31.110,43.400,19.050,51.110,68.750,53.570,30.000,56.520)
\tkzKiviatLine[very thick,color=color1](23.910,23.910,37.780,37.740,40.480,68.890,59.380,60.710,17.500,65.220)
\tkzKiviatLine[very thick,color=color2](39.130,43.480,42.220,41.510,23.810,48.890,40.620,64.290,30.000,60.870)
\tkzKiviatLine[dashed,very thick,color=color9](77.780,40.000,66.670,0,0,0,0,0,59.090,79.310)

\tkzKiviatGrad[prefix=,unity=20,suffix=\%](1)
\node[align=center] at (-5.5,5.5) {\normalsize\textbf{Forest}};
\node[align=center] at (5.5,5.5) {\normalsize\textbf{City}};
\draw[red,thick,dashed] (0,-6) -- (0,6);
\end{tikzpicture}
\caption{Proprietary models}
\end{subfigure}
\begin{subfigure}{0.47\linewidth}
\centering
\begin{tikzpicture}
\tkzKiviatDiagram[scale=0.35,
        label distance=.5cm,
        radial  = 10,
        gap     = 1,
        lattice = 5,
        step=.05]
        {Trash\\pile,Fire,Car,Trash\\pile,Person,Building,Campsite,Forest\\fire,Road\\construction,Crowd}
\tkzKiviatLine[very thick,color=color6](13.040,21.740,17.780,9.430,4.760,28.890,6.250,35.710,12.500,30.430)
\tkzKiviatLine[very thick,color=color7](2.170,30.430,4.440,9.430,9.520,15.560,25.000,42.860,12.500,4.350)
\tkzKiviatLine[very thick,color=color8](21.740,21.740,20.000,26.420,16.670,62.220,43.750,46.430,10.000,43.480)
\tkzKiviatLine[dashed,very thick,color=color9](77.780,40.000,66.670,0,0,0,0,0,59.090,79.310)

\tkzKiviatGrad[prefix=,unity=20,suffix=\%](1)
\node[align=center] at (-5.5,5.5) {\normalsize\textbf{Forest}};
\node[align=center] at (5.5,5.5) {\normalsize\textbf{City}};
\draw[red,thick,dashed] (0,-6) -- (0,6);
\end{tikzpicture}
\caption{Open source models}
\end{subfigure}

\vspace{0.5cm}

\begin{tikzpicture} 
    \begin{axis}[%
    hide axis,
    xmin=10,
    xmax=50,
    ymin=0,
    ymax=0.4,
    legend columns=4,
    legend style={draw,legend cell align=left,column sep=1ex}
    ]
\addlegendimage{color0,ultra thick}
\addlegendentry{GPT-4o};
\addlegendimage{color1,ultra thick}
\addlegendentry{Claude 3.5 Sonnet};
\addlegendimage{color2,ultra thick}
\addlegendentry{Gemini 2.0 Flash};
\addlegendimage{color6,ultra thick}
\addlegendentry{Qwen2-VL-72B};
\addlegendimage{color7,ultra thick}
\addlegendentry{Llava-Onevision 72b};
\addlegendimage{color8,ultra thick}
\addlegendentry{Pixtral-Large 124B};
\addlegendimage{color9,ultra thick,dashed}
\addlegendentry{Human baseline};
\end{axis}
\end{tikzpicture}
\end{scriptsize}
\caption{\textbf{Success rate per class:} In this figure, we show the model success rate per each target class for both the Forest and the City environments. Classes with large visible objects such as \textit{Building} for Forest or \textit{Crowd} for City are significantly easier for most models to locate. On the other hand, classes that are difficult to distinguish from the background, such as both garbage classes, are only located by more advanced models. Human baseline for City is marked with dashed line.}
\label{fig:accuracy-per-class}
\end{figure*}